\documentclass[12pt]{article}
\usepackage[utf8]{inputenc}
\usepackage[T1]{fontenc}
\usepackage{hyperref}
\title{Yamba Language Resources for K-12 Curriculum Development}
\author{Compiled by Claude Sonnet 5.5 (claude-sonnet-5-5) for Andr\'as Kornai}
\date{}
\begin{document}
\maketitle

\section{Introduction}
Yamba (ISO 639-3 \texttt{yam}) is a Grassfields language of the Northwest Region of Cameroon, with a few speakers in Nigeria, and about 80,000 speakers. It uses the General Alphabet for Cameroon Languages. The search found a provisional online dictionary, scripture and audio resources, and an older grammar.

\section{Primary Educational Resources}
No primers or textbooks were found in the search.

\section{Linguistic and Reference Materials}
\begin{itemize}
\item Yamba--English provisional dictionary on Webonary, with more than 3,300 entries based on the Mbem dialect: \url{https://www.webonary.work/yamba/overview/introduction/}
\item Webonary Yamba site: \url{https://www.webonary.work/yamba/?p=9444}
\item Barnwell (1969), a reference listed in Glottolog. Glottolog files it under Mbembe Cross River, so check that it is about this language: \url{https://glottolog.org/resource/reference/id/105907}
\item Wikipedia, Yamba language: \url{https://en.wikipedia.org/wiki/Yamba_language}
\item Omniglot, Yamba (writing system). It conflicts with other sources on where Mbem is located: \url{https://omniglot.com/writing/yamba.htm}
\end{itemize}

\section{Religious and Community Resources}
\begin{itemize}
\item Scripture Earth, Yamba. It lists the Gospel of Mark audio, the JESUS Film and the LUKE Video: \url{https://www.scriptureearth.org/00eng.php?iso=yam}
\end{itemize}

\section{Digital and Audio Resources}
\begin{itemize}
\item Global Recordings Network, Yamba (Nkot) audio: \url{https://globalrecordings.net/en/language/6269}
\item Wikitongues page for \texttt{yam}: \url{https://wikitongues.org/languages/yam/}
\end{itemize}

\section{Curriculum Development Context}
A dictionary and an alphabet exist, which gives a base for materials. No school programme or textbooks were found in the search. Omniglot mentions radio programmes, which I did not verify. A curriculum would build on the Webonary dictionary and the Cameroon alphabet, with community and local language-committee input.
\end{document}